\section{Limitations and Future Work}
\label{sec:limitations}

\textsc{ShapeCodeBench} makes deliberate scoping choices in its v1 incarnation, each
of which is a direction for future work.

\paragraph{Monochrome palette.}
V1 is black-on-white. This collapses draw-order sensitivity: later shapes can
paint foreground pixels but cannot erase or overwrite earlier ones. A richer
palette (multiple colors or an explicit \texttt{clear} primitive) would make
draw order first-class and enable sharper compositionality tests.

\paragraph{Four primitives.}
The DSL currently covers only filled and hollow circles and squares. Extending
to rectangles, lines, polygons, or parametric curves would stress different
kinds of visual reasoning and likely move saturation further out.

\paragraph{Zero-shot only.}
We evaluate without chain-of-thought prompting or few-shot examples. These are
natural knobs to explore and may change the ordering of models, especially for
the reasoning-heavy hard tier.

\paragraph{Model-inference variability.}
The frozen evaluation images, parser, renderer, and scorer are deterministic:
regenerating \textsc{eval\_v1} from the published seeds should reproduce the
same target PNGs and metric computation. The remaining variability is in model
inference. Repeating the same image-and-prompt request to a hosted multimodal
model may produce a different predicted program, so we cannot guarantee
bit-exact reproduction of the reported model scores. Each run's
\texttt{run\_config.json} records the model configuration and invocation settings
needed to reproduce the experimental setup.

\paragraph{No human baseline.}
We do not report human performance. An informal human baseline on a small
sample of scenes would make the ``how hard is this really?'' claim concrete,
and we plan to add it in a subsequent revision.

\paragraph{Model coverage.}
We evaluate two frontier multimodal models -- \textsc{Claude Opus 4.7} (1M
context) and \textsc{GPT-5.5} -- across two reasoning-effort tiers each. Other
frontier multimodal systems are not covered here but can be added by
implementing a new adapter against the same \texttt{ModelAdapter} Protocol.

\paragraph{Contamination resistance.}
The current eval split's seeds are public. We rely on the \emph{renewability}
of the benchmark rather than on secrecy: to evaluate under a contamination-free
setting, regenerate \textsc{eval\_vN} from fresh seeds using the published
generator and manifest. This is intentionally cheap: the entire 150-sample split
regenerates in under a second.

\paragraph{Training signal, not a training pipeline.}
The current evaluator is an offline benchmark, not a differentiable pretraining
loss. Predictions are discrete DSL text, the parser is symbolic, the renderer is
Pillow-based, and the reported metrics are computed after hard rasterization.
Future training use would therefore require a separate setup, such as supervised
fine-tuning on generated image/program pairs, reinforcement learning or
rejection sampling with render-based rewards, a learned reward or critic model,
or a differentiable renderer variant. Any such use should train on generated
training seeds and reserve \textsc{eval\_v1} or fresh held-out splits for clean
evaluation.

\paragraph{Reproducibility.}
We publish the benchmark code, the frozen \textsc{eval\_v1} dataset with
SHA-256 hashes per image, per-run artifacts, analysis scripts, and paper sources.
\texttt{docs/REPRODUCIBILITY.md} walks through the end-to-end workflow from a
clean checkout.

\section*{Acknowledgments}

We thank the authors of the prior benchmarks cited throughout this paper for
laying the groundwork that \textsc{ShapeCodeBench} builds on, and the maintainers of
\texttt{Pillow}, \texttt{numpy}, and \texttt{scipy} for tools that made the
evaluator and heuristic baseline straightforward to implement.
